\documentclass[
 reprint, %preprint for single col, double spacing
 %pre, % others journal styles: pra, prl, etc
 amsmath, amssymb, aps, % preload maths packages
 a4paper,
 %onecolumn
]{revtex4-2}

\usepackage{microtype}
\usepackage{bm}
\usepackage{lipsum} % just to fill text
\usepackage{dcolumn} % for decimal alignment in tables
\usepackage{siunitx}
\usepackage{physics}

\usepackage{graphicx}
% be careful here. I often get errors when using packages like 
% subfig, subfigure, subcaption
% they often break or give weird warnings when using with hyperref
%\usepackage{subfigure}
\usepackage{subfig}

\usepackage{natbib} % this is the package I recommend using or referencing and citations
% put your references in a `.bib' file (see ref.bib) 
% this is loaded at the end using the \bibliography{<name>} command
% how the bibliography looks depends on the style. We can change from the default with the command: 
\bibliographystyle{unsrtnat} % this is probably a good one to go with
%\bibliographystyle{unsrt}
%\bibliographystyle{siam}
%\bibliographystyle{apsrev4-2}
% many styles exist and the style you should use will depend on what is expected. 
% Imagine they want:
% numbered, square brackets for the citations
% name of authors, article title, journal name, and DOI if available

% revtex already loads hyperref but explicit import is needed to make cleveref work
\usepackage[hidelinks]{hyperref}
\usepackage{cleveref}

\begin{document}

\title{RevTex example}

\author{Jeremy Worsfold}
\email{j.worsfold@unsw.edu.au}
\affiliation{
    School of Physics,
    University of New South Wales - Sydney 2052, Australia
}

\date{\today}

\begin{abstract}
\lipsum[1]

\end{abstract}


\maketitle


\section{Introduction}
\label{sec:intro}

\lipsum

\section{Tables and figures}

You can force tables and figures to be in specific areas using the \texttt{[h]} option but generally, I think it's best to let the compiler find where it should go.
Other options include \texttt{[t]} for top of the page, or \texttt{[b]} for bottom of the page. Usually in journals like APS, figures go at the top.
You can also just move it about in the text to move the figure/table onto another page. But make sure it's close or directly after being referred to in the text.

\begin{table}[h]
    \centering
    \begin{tabular}{cccc}
         Col 1 (\unit{\newton\meter})& Col 2  & Col 3  & Col 4 \\
         \colrule
         1 & 2 & 3 & 4 \\
         5 & 6 & 7 & 8 \\
    \end{tabular}
    \caption{Caption}
    \label{tab:my_label}
\end{table}

\begin{table}[h]
    \begin{ruledtabular}
    \begin{tabular}{lcdr}
        Left & Centered & \multicolumn{1}{c}{Decimal} & Right\\
        \colrule
        1 & 2 & 3.001 & 4 \\
        10 & 20 & 30 & 40 \\
        100 & 200 & 300.0 & 400 \\
    \end{tabular}
    \end{ruledtabular}
    \caption{An example table}
    \label{tab:table1}
\end{table}


\section{Other Packages}


\subsection{SI units}
%\usepackage{siunitx}
An example use of units $\SI{52.3}{\newton}=\SI{52.3}{\kg\per\meter\squared}$.

\subsection{physics package}
\newpage
\begin{equation}
    \int_{-\infty}^\infty e^{-x^2} d x = \sqrt{\pi} \label{eq:gaussInt}
\end{equation}

\begin{equation}
    \int_{-\infty}^\infty e^{-x^2} \dd x = \sqrt{\pi}
\end{equation}

\begin{equation}
    \frac{\partial^2u}{\partial x^2} = \pdv[2]{u}{x} \label{eq:partialDiffExample}
\end{equation}

\newcommand{\myvariable}{\gamma}

\begin{align}
    \myvariable
\end{align}

\begin{equation}
    \frac{d ^2u}{\partial x^2} = \pdv[2]{u}{x} \label{eq:partialDiffExample}
\end{equation}

\section{referencing text}
\label{sec:ref}

This is a reference to Section \ref{sec:intro}. This is a neater reference to \cref{sec:intro} or \Cref{sec:intro}.
It also works for equations such as \cref{eq:gaussInt}, as opposed to Eq. \eqref{eq:gaussInt}.
However, if you directly mention that the thing you're referring to is an equation, ODE, PDE etc, then it's probably better practice to just give the equation number. 
For example, substituting \cref{eq:gaussInt} into the PDE \eqref{eq:gaussInt} and simplifying, we obtain...
You can also reference figures like \cref{fig:doublefig} or by sub figure: \cref{subfig:a}.

\subsection{Citations}

I recommend you use \texttt{natbib} to do referencing, this works well with the \texttt{revtex} documentclass which also comes with a bibliography style. 
In this `vancouver' referencing style, you only refer using a number so the only command you really need is \texttt{cite}. For instance, this is a nod to a great mathematician \cite{turing1990chemical}. 
Citing multiple sources at the same time just requires separating each key with commas \cite{turing1990chemical,cooley1967fft}.

Give your entries in the bibliography useful keys and try to keep to a convention, e.g. \texttt{<lead author><date><first word of title>}.
Reference managers (and google scholar) will automatically do things like this. Wrapping text with \{ \} brackets ensures words stay capitalized.


\begin{figure*} % make it full page width
    \centering
    \subfloat{\includegraphics[width=8.6cm]{figures/fig_1.eps}}
    \hfill
    \subfloat{\includegraphics[width=8.6cm]{figures/fig_1.eps}}
    \caption{An example of sub-figures}
    \label{fig:doublefig}
\end{figure*}

\begin{figure*}
    \centering
    \subfloat[\label{subfig:a}]{%
        \includegraphics[width=0.3\linewidth]{figures/fig_1-eps-converted-to.pdf}%
    }
    \hfil
    \subfloat[\label{subfig:b}]{%
        \includegraphics[width=0.3\linewidth]{figures/fig_1-eps-converted-to.pdf}%
    }
    \caption{A graphical depiction of the setup}
    \label{fig:subfigureExample}
\end{figure*}

% should have the bibliography before the appendices generally.
\bibliography{ref}

\appendix

\section{First appendix}

The numbering of equations and sections changes after the appendix
\begin{equation}
    e^{i\pi}\equiv-1.
\end{equation}

\end{document}